\documentclass[manuscript,screen]{acmart}

% Remove ACM-specific copyright and conference metadata placeholders since none were supplied
\settopmatter{printacmref=false}
\renewcommand\footnotetextcopyrightpermission[1]{}
\pagestyle{plain}

\usepackage{listings}
\usepackage{xcolor}

\lstset{
  basicstyle=\ttfamily\footnotesize,
  keywordstyle=\color{blue!70!black}\bfseries,
  commentstyle=\color{gray}\itshape,
  stringstyle=\color{teal!70!black},
  breaklines=true,
  frame=single,
  framesep=4pt,
  rulecolor=\color{black!20},
  showstringspaces=false,
  tabsize=2,
  morekeywords={var, val, by, remember, mutableStateOf, LaunchedEffect, while, true, delay}
}

\begin{document}

\title{EchoBox: A Jetpack Compose Android Application for Temporal Self-Communication and Memory Preservation}

\author{The EchoBox Development Team}
\affiliation{%
  \institution{EchoBox Open-Source Project}
  \city{Global}
  \country{}
}

\begin{abstract}
As digital footprints grow increasingly ephemeral and real-time oriented, tools for intentional, delayed self-reflection remain scarce. This paper introduces EchoBox, an open-source Android application developed using modern native technologies---specifically Jetpack Compose, Kotlin Coroutines, and Room Database---designed to facilitate temporal self-communication. EchoBox enables users to compose text messages, attach multimedia, assign categorical tags, and logically time-lock these entries until a specified future date. A live, second-by-second ticker ensures real-time state updates upon maturity. By bridging personal journaling with temporal constraints, EchoBox provides a structured, privacy-focused medium for long-term goal setting, psychological grounding, and asynchronous dialogue with the future self.
\end{abstract}

\maketitle

\section{Introduction}
Modern mobile applications are heavily optimized for instant gratification, high-frequency synchronization, and immediate social feedback. Conversely, psychological principles highlight the cognitive and emotional benefits of delayed reflection---specifically, communicating aspirations, emotional states, or advice to one's future self.

To explore this paradigm in a mobile context, we developed EchoBox, a native Android application that bridges personal journaling with strict temporal constraints. EchoBox allows users to create ``time capsules''---sealed packages containing text, metadata tags, and optional photographs---that remain inaccessible until a pre-determined future date and time are reached.

\section{System Architecture \& Tech Stack}
EchoBox is engineered following modern Android application architecture guidelines, prioritizing reactive UI patterns and robust local persistence:
\begin{itemize}
    \item \textbf{User Interface:} Built entirely with Jetpack Compose, ensuring a fluid, declarative layout and responsive state handling.
    \item \textbf{Architecture \& Concurrency:} State management is driven by Kotlin StateFlow, and asynchronous background tasks (such as disk I/O) are managed safely via Kotlin Coroutines (\texttt{Dispatchers.IO}).
    \item \textbf{Local Persistence:} Utilizes the Room Database library (\texttt{MessagesDao}, \texttt{MessagesEntity}) for reliable, structured local storage of message metadata.
    \item \textbf{Media Management:} Integrates the native Android Photo Picker (\texttt{PickVisualMedia}) and Coil for efficient, memory-safe asynchronous image rendering.
\end{itemize}

\section{Core Implementation Mechanics}
The core innovation of EchoBox lies in its real-time conditional unlocking mechanism and localized storage pipeline.

\subsection{Real-Time State Evaluation}
Unlike server-dependent architectures that rely on push notifications or cron jobs to unlock content, EchoBox implements a lightweight reactive ticker inside the \texttt{MainScreen} composable:

\begin{lstlisting}[language=Java]
var currentTime by remember {
    mutableStateOf(System.currentTimeMillis())
}
LaunchedEffect(Unit) {
    while (true) {
        delay(1000L) // Refresh every second
        currentTime = System.currentTimeMillis()
    }
}
\end{lstlisting}

This ticker continuously re-evaluates whether a message's \texttt{targetedDate} has passed. When \texttt{targetedDate <= currentTime}, the UI state transitions automatically, prompting the user with a distinct visual indicator and activating the ``Open Capsule'' action button.

\subsection{Image Storage Pipeline}
Attached images are handled securely to prevent memory leaks and redundant file duplication. Selected URIs are processed through an \texttt{ImageStorageManager} that copies assets into app-private internal storage, persisting file paths inside the Room database while ensuring orphaned images are properly garbage-collected upon deletion.

\section{Discussion and Use Cases}
EchoBox serves multiple psychological and productivity-oriented use cases:
\begin{enumerate}
    \item \textbf{Goal Tracking \& Accountability:} Users can set 6-month or 1-year goals, sealing their intentions and reviewing whether milestones were achieved upon unlocking.
    \item \textbf{Emotional Reflection:} Writing letters during periods of high stress or joy to be read during future reflective periods.
    \item \textbf{Privacy-First Design:} Because all data and images are stored locally using Room and app-private storage, users retain absolute ownership of their historical records without external cloud exposure.
\end{enumerate}

\section{Conclusion and Future Work}
EchoBox demonstrates how modern Android frameworks (Jetpack Compose, Room, Coroutines) can be harmonized to build a performant, privacy-respecting temporal journaling tool. Future iterations of EchoBox will explore optional, zero-knowledge end-to-end encrypted cloud synchronization, and AI-driven sentiment analysis trends over multi-year archives.

\end{document}